% Root isolated candidate; full independent audit and finite supplement pending.
\subsection{Exact symmetric image in every odd primitive degree}

\begin{theorem}[The entire elliptic double-cover family]
\label{thm:root-all-odd-double-cover-symmetric-image}
\status{Candidate; complete written proof awaiting independent audit and supplement.}
Let \(E:y^2+y=x^3\) over \(\mathbf F_2\), let
\(A=E\times E\), and put
\[
 L_m=\mathcal O_A(mV+H+\Gamma),\quad
 V=\{O\}\times E,\quad H=E\times\{O\},\quad
 \Gamma=\{(x,y):x+y=0\},\quad n=2m+1,\quad m\ge1.
\]
On a nonempty parameter open the curves in \(|L_m|\) are
smooth and connected, have genus \(n+1\), and have
\(2\)-rank \(n-1\). Their coherent relative Frobenius has
generic stable rank \(n-1\), equal to its full rank.
The actual geometric and arithmetic relative Lang images are
both the faithful symmetric group
\[
 S_n\longrightarrow\operatorname{GL}
       (\mathbf F_2^n/\mathbf F_2\mathbf1).
\]
The determinant character is trivial. For \(n=3\) this is
\(\operatorname{GL}_2(\mathbf F_2)\); for every odd
\(n\ge5\) it is a proper subgroup of
\(\operatorname{GL}_{n-1}(\mathbf F_2)\).
\end{theorem}
\begin{proof}
Let \(E:y^2+y=x^3\) be supersingular over an algebraically closed
field of characteristic two. Put \(A=E\times E\), let \(q:A\to E\)
be the first projection, and use
\[
 V=\{O\}\times E,\quad H=E\times\{O\},\quad
 \Gamma=\{(x,y):x+y=0\},\quad
 L_0=\mathcal O_A(H+\Gamma),\quad L_m=L_0(mV).
\]
The already audited product-line lemma proves, for every \(m\ge1\),
that \(L_m\) is ample and globally generated, of type \((1,n)\)
with \(n=2m+1\), and that
\(\mathcal V_0=q_*L_0\) is a rank-two degree-one bundle.
In particular \(h^0(A,L_m)=n\). This input does not require a
prime degree or any generic Cartier-rank theorem.

\paragraph{The actual ruled quotient.}
On each fibre of \(q\), the restriction of \(L_0\) is the complete
degree-two line bundle \(\mathcal O_E(O+(-x))\). Fibrewise evaluation
therefore gives a morphism
\[
 \pi:A\longrightarrow S=\mathbf P_E(\mathcal V_0),\qquad
 \pi^*\mathcal O_S(1)=L_0,
\]
where the projective-bundle convention parametrizes one-dimensional
quotients. The map is proper and fibrewise finite of degree two, hence
finite by the proper finite-fibre criterion of
\cite[Lemma~37.44.1, Tag~02LS]{StacksProperFiniteCP38}.
The degree-two elliptic linear-system involution is
\[
                 \iota(x,y)=(x,-x-y).
\]
Indeed two points on the fibre over \(x\) define the same
degree-two divisor precisely when their sum is \(-x\). Thus
\(\pi\iota=\pi\); its degree-two function-field extension is
separable and has this nonidentity automorphism. On an affine
open of the smooth, hence normal, surface \(S\), the invariant
subring of the finite coordinate algebra of its inverse image
is integral over the base ring and contained in its fraction
field. Normality makes this invariant subring equal to the base
ring. Hence \(\pi\) is the actual quotient. The displayed
pullback of \(\mathcal O_S(1)\) supplies descent of the line
bundle itself, not just invariance of its numerical class.

\paragraph{The fixed elliptic curve and its quotient.}
The fixed scheme is the smooth elliptic curve
\[
                   T=\{(x,y):x=-[2]y\}\simeq E.
\]
Its tangent equation is \(dx=0\), since the differential of \([2]\)
vanishes. The differential of \(\iota-1\) has rank one at every
point of \(T\). Locally at a fixed point choose a parameter \(v\)
whose differential is nonzero in the vertical direction, and set
\(u=\iota(v)+v\). Then \(u,v\) are regular parameters,
\(\iota(u)=u\), \(\iota(v)=v+u\), and \(T=(u=0)\).
Invariants commute with this completion: they are the kernel
of the endomorphism \(\iota-1\) of the finite pushforward
module, and completion of the noetherian base is flat. A fixed
point is the only point above its quotient image, since its
degree-two fibre divisor is twice that point. Thus the completed
invariant ring is exactly
\[
 k[[u,v]]^{\langle\iota\rangle}
                  =k[[u,w]],\qquad w=v^2+uv.
\]
Indeed \(k[[u,v]]\) is free of rank two over the displayed
subring, with basis \(1,v\); invariance of \(a+vb\) forces
\(ub=0\), hence \(b=0\). Consequently the reduced image
\(R=\pi(T)\) is smooth, and \(T\to R\) is a purely inseparable
degree-two map, locally \(w=v^2\). Thus \(R\) has genus one.
The original map \(q|_T=[-2]\) is purely inseparable of degree four,
because \(E\) is supersingular. Hence the induced \(R\to E\)
is purely inseparable of degree two as well. No geometric branch
points are identified by either map.

\paragraph{Restriction is a complete degree-\(n\) system.}
Set
\[
 \mathcal D_m=\mathcal O_S(1)\otimes\bar q^*\mathcal O_E(mO),
 \qquad N_m=\mathcal D_m|_R,
\]
where \(\bar q:S\to E\). All sections of \(L_m\) descend to
\(\mathcal D_m\): the pullback map on sections is an isomorphism,
as follows from the relative evaluation construction and
\(\bar q_*\mathcal D_m=\mathcal V_0(mO)=q_*L_m\).
Intersection on \(A\) gives
\[
 T\cdot V=4,\qquad T\cdot H=1,\qquad T\cdot\Gamma=1,
 \qquad L_m\cdot T=4m+2=2n.
\]
Since \(T\to R\) has degree two, \(\deg N_m=n\).
The restriction
\[
          H^0(S,\mathcal D_m)\longrightarrow H^0(R,N_m)
\]
is injective. A section in its kernel would pull back to a
nonzero section of \(L_m\) vanishing on \(T\), whereas
\((L_m-T)\cdot V=2-4=-2\) and \(V\) is nef, so
\(H^0(A,L_m-T)=0\). Both spaces have dimension \(n\), the
second by Riemann--Roch on the elliptic curve. Restriction is
therefore an actual linear isomorphism. Equivalently, restriction
to \(T\) consists of the pullbacks along \(T\to R\) of the
complete degree-\(n\) system, so locally these restrictions are
Frobenius squares. This conclusion uses an actual descended
bundle and quotient, rather than choosing arbitrary square roots
of individual sections.

\paragraph{The generic branch set and rank.}
For generic \(s\), the divisor \(D_s\) of \(\mathcal D_m\)
is a section of \(\bar q\). A section of the globally generated
rank-two bundle \(\mathcal V_0(mO)\) is generically nowhere zero:
the incidence of a zero at some point of the curve has codimension
at least one in its section space. Thus \(D_s\simeq E\).
Its pullback \(C_s\in|L_m|\) is smooth and connected on a
nonempty open by the already proved odd globally generated
smooth-member lemma. Independently, the complete degree-\(n\)
system on \(R\) has a nonempty open of reduced divisors, so
\(D_s\cap R\) consists of \(n\) distinct points generically.
These opens meet because the coefficient space is integral.

Off \(R\), the quotient \(\pi\) is etale. One can
see this without a tame-quotient assertion: at an orbit of two
distinct points, the completed finite algebra is the product of
their completed local rings. The involution exchanges the two
factors, and its invariants are one factor. The completed local
map at either point is consequently an isomorphism. On \(R\), its
local equation is \(v^2+uv=w\). Thus the separable double cover
\(C_s\to D_s\simeq E\) has exactly the geometric branch
points \(D_s\cap R\). It is separable because \(D_s\) is not
contained in \(R\), so its quadratic equation has nonzero linear
coefficient in the generic field. The purely inseparable map
\(R\to E\) preserves the geometric branch permutation set,
with its Galois action. The canonical double-cover sheaf proof
already audited for the prime examples now gives
\[
 H^1_{\mathrm{et}}(C_{\bar\eta},\mathbf F_2)
          \simeq \mathbf F_2^{\mathcal B}/\mathbf F_2,
 \qquad |\mathcal B|=n.
\]
In particular the generic \(2\)-rank is \(n-1\), without invoking
the finite list of certified primes. Adjunction gives genus
\(n+1\), as before.

\paragraph{Exact symmetric branch monodromy.}
Under the restriction isomorphism, the branch family is the
universal reduced divisor of \(|N_m|\) on \(R\). Its ordered
pair incidence is a projective bundle over
\(R\times R\setminus\Delta\): two distinct evaluations are
independent since \(\deg N_m(-P-Q)=n-2>0\). This incidence
is irreducible, so the generic branch permutation group is
two-transitive.

Choose a divisor with one exact double point and every other
point simple. Such a divisor exists: for \(n\ge5\), the
residual degree \(n-2\) system is globally generated and its
general divisor is reduced away from the chosen point; for
\(n=3\), choose a point outside the finite set
\(N_m\simeq\mathcal O_R(3P)\). The latter set is finite since
\([3]\) is finite etale in characteristic two.

The local universal polynomial at the double point has, after
formal Weierstrass preparation, the form
\[
                       X^2+aX+b.
\]
The first two evaluation jets are independent, so \(a,b\)
are independent formal coefficient parameters. In detail, at
the chosen exact double point the quadratic leading coefficient
is nonzero, while the differential of \(s\mapsto(b,a)\)
is an invertible triangular transform of evaluation of \(s\)
and its first derivative there. The map to these two jets is
surjective because \(\deg N_m(-2P)=n-2>0\).
The other \(n-2\) simple roots lift uniquely over the complete
local parameter ring. Localizing that ring at \((a)\) and
strict-henselizing gives a trait whose generic point maps to the
reduced-divisor open; its finite permutation image is therefore
a subgroup of the global geometric image. At the reduced
divisor \(a=0\), the element \(b\) is nonsquare in the
strict-henselian residue field; separably closed does not imply
perfect here. Adjoining a root gives a separable quadratic
extension of the generic field, with automorphism
\(X\mapsto X+a\) and a purely inseparable residue extension.
Its full inertia is therefore \(C_2\). All other roots are
etale and fixed after strict henselization. Nonsquareness of
the residue of \(b\) follows from \(db\ne0\) in the
residue fraction field of the power-series ring; separable
algebraic extension preserves this nonzero differential. Thus
the quadratic factor is irreducible there and its order-two
automorphism actually occurs in inertia. This supplies an
isolated transposition by fierce characteristic-two inertia;
it is not a tame-tangency argument. Two-transitivity and a
transposition imply the full group \(S_n\).

Finally its action on \(\mathbf F_2^n/\mathbf F_2\) is faithful
for odd \(n\ge3\): if a permutation acts trivially, each
\(e_i+e_{g(i)}\) is constant, but a nonzero such vector has
weight two, which is neither zero nor \(n\). Thus every index
is fixed. This identifies the actual relative geometric
Lang image with \(S_n\). The canonical branch-module
isomorphism is Galois-equivariant for the arithmetic action as
well. All the objects were defined over \(\mathbf F_2\),
so the arithmetic image is contained in this same \(S_n\)
and contains the full geometric image; it is equal to it.
The determinant is trivial because \(\mathbf F_2^\times=1\).
The ambient surface has zero stable cohomology, so the natural
Artin--Schreier and stable hypersurface sequences identify this
module with the requested relative Lang module, exactly as in
Lemma~\ref{lem:root-double-cover-branch-module} and
Theorem~\ref{thm:root-primitive-prime-permutation-counterexamples}.
No prime-dependent Cartier certificate has been used.


The stated properness for every \(n\ge5\) follows from the
integer order comparison already proved for all integers
\(n\ge5\): it starts with \(120<20160\), and the next
linear-group order grows by \(2^{n-1}(2^n-1)>n+1\).
\end{proof}

\begin{corollary}[Irreducibility and the exact connected components]
\label{cor:root-all-odd-double-cover-orbits}
\status{Candidate; complete written proof awaiting independent audit and supplement.}
For every odd \(n=2m+1\ge3\), the relative
\(\mathbf F_2[S_n]\)-module is irreducible. The nonzero
relative Lang cover has exactly \((n-1)/2\) geometric
connected components, of degrees
\[
              \binom n1,\binom n2,\ldots,
                             \binom n{(n-1)/2}.
\]
Each component is defined over \(\mathbf F_2\).
\end{corollary}
\begin{proof}
Since \(n\) is odd, projection modulo the all-one line
identifies the quotient with the even-sum subspace of
\(\mathbf F_2^n\). Any nonzero vector there has two
coordinates with different entries. Subtracting its image under
the transposition of those coordinates gives \(e_i+e_j\).
Thus any nonzero invariant subspace contains one such vector,
and conjugating under \(S_n\) gives every \(e_i+e_j\).
These span the even-sum subspace, proving irreducibility.

Every nonzero class has a unique subset representative of
cardinality \(1\le j\le(n-1)/2\), up to taking its
complement. The full symmetric group is transitive on each
set of a fixed cardinality. These are exactly the nonzero
orbits and have the displayed binomial sizes. Finite-cover
components correspond to these orbits. Arithmetic and geometric
images agree and preserve every cardinality, giving the final
descent assertion.
\end{proof}

\begin{remark}[Scope and reproducible evidence]
\label{rem:root-all-odd-double-cover-scope}
\status{Candidate scope; no fixed-embedding large-degree conclusion.}
The proof extends the actual prime examples to every odd degree,
and determines their exact image and components. The fibre
degree remains two, so none of these lines is very ample and
this varying primitive family is not the sequence of powers of
a fixed embedding line. General A.4 remains open.
The original proper finite-fibre criterion and its complete
proof were opened at Tag~02LS; the noetherian formal-functions
proof at Tag~02OG was also read in full. The exact quotient,
formal invariant ring and wild quadratic inertia are proved
above. A separate symbolic supplement is being prepared; this
candidate is not promoted before that program/output and the
complete independent written audits are retained.
\end{remark}
